\chapter{曲线积分和曲面积分}
在一元函数的积分中，通过计算$f(x)$的原函数$F(x)$满足：
\[\dd{F(x)} = f(x)\dd{x}\]
进一步基于Newton-Leibniz公式通过区域的边界计算区域内部导数积分的值.而在计算重积分的过程中，往往是使用各种方法将其转换为累次积分，再使用Newton-Leibniz公式.

是否存在、能否求取多元函数的原函数$F(x,y)$使得：
\[\dd{F(x,y)} = F'_x(x,y)\dd{x} + F'_y(x,y)\dd{y} = \nabla F \cdot \dd{\vb{r}}\]
进而推广出高维空间的Newton-Leibniz公式，将区域内部导数的积分与区域边界上的函数值相联系？这就是Gauss-Stokes公式.


\section{数量场在曲线上的积分}
要想求取原函数$F(x,y)$，就要尝试作积分：
\[\iint_D \dd{F(x,y)} = \int_L \nabla F \cdot \dd{\vb{r}}\]
显然，后者是一个向量场上的积分，在解决这个问题之前，首先要解决数量场上的积分.